%%%%%%%%%%%%%%%%%%%%%%%%%%%%%%%%%%%%%%%%%%%%%%%%%%
\begin{frame}[fragile]\frametitle{Reading the Gates of an LSTM}
\begin{itemize}
\item The memory $c_t$ is a running tally; three gates control it
\item Each gate is a number between 0 and 1 for every memory slot
\item Forget gate $f_t$: how much of the old memory to keep
\item Input gate $i_t$: how much of the new candidate $\tilde{c}_t$ to write
\item Output gate $o_t$: how much of the memory to show as the output $h_t$
\item $c_t = f_t c_{t-1} + i_t \tilde{c}_t$ keeps part of the old and adds part of the new
\item A GRU does the same job with two gates and no separate memory
\end{itemize}
\end{frame}

%%%%%%%%%%%%%%%%%%%%%%%%%%%%%%%%%%%%%%%%%%%%%%%%%%
\begin{frame}[fragile]\frametitle{Quick Check: RNN Concepts}
\begin{block}{Think About It}
1. What does the hidden state of an RNN carry from one time step to the next?
\begin{enumerate}
  \item[A.] Only the previous input.
  \item[B.] A running summary of the sequence seen so far.
  \item[C.] The loss value.
  \item[D.] The learning rate.
\end{enumerate}
2. The same weights are used at every time step. What happens to the number of parameters when the sequence gets longer?
\begin{enumerate}
  \item[A.] It grows with the length.
  \item[B.] It stays the same.
  \item[C.] It doubles at every step.
  \item[D.] It shrinks.
\end{enumerate}
\end{block}
\end{frame}

%%%%%%%%%%%%%%%%%%%%%%%%%%%%%%%%%%%%%%%%%%%%%%%%%%
\begin{frame}[fragile]\frametitle{Quick Check: RNN Concepts}
\begin{block}{Answer}
\textbf{1. Correct: B.} The new state mixes the current input with the previous state, so it summarizes everything seen so far. A would describe a network that only looks back one input.

\textbf{2. Correct: B.} The same weights are reused at every step, so a longer sequence means more steps through the same weights, not more weights. This is why the model can handle sequences of any length.
\end{block}
\end{frame}

%%%%%%%%%%%%%%%%%%%%%%%%%%%%%%%%%%%%%%%%%%%%%%%%%%
\begin{frame}[fragile]\frametitle{}
\begin{center}
{\Large Case Study: Forecasting a Sensor Signal}
\end{center}
\end{frame}

%%%%%%%%%%%%%%%%%%%%%%%%%%%%%%%%%%%%%%%%%%%%%%%%%%%%%%%%%%
\begin{frame}[fragile]\frametitle{The Data}
\begin{itemize}
\item 3,650 daily minimum temperatures in degrees C, Melbourne, 1 January 1981 to 31 December 1990
\item Not machine data, but it has the shape of sensor data: regular samples, a slow yearly swing, and day-to-day noise on top
\item Values run from 0.0 to 26.3 with mean 11.18 and standard deviation 4.07
\item 29 February of 1984 and 1988 is absent from the file; we treat the series as consecutive days
\end{itemize}
\begin{lstlisting}
import pandas as pd

df = pd.read_csv('daily-minimum-temperatures.csv')
series = df['Temperature'].to_numpy('float32')

>>> len(series), series.min(), series.max()
(3650, 0.0, 26.3)
\end{lstlisting}
\end{frame}

%%%%%%%%%%%%%%%%%%%%%%%%%%%%%%%%%%%%%%%%%%%%%%%%%%%%%%%%%%
\begin{frame}[fragile]\frametitle{From a Series to Windows}
\begin{itemize}
\item The task: given the last 7 days, predict tomorrow
\item Slide a window of 7 days along the series; each window is one example, the day after it is the answer
\item Split by time, never at random: the first 8 years (1981 to 1988) train, the last 2 (1989 and 1990) test
\item Scale with the training period's mean (11.11) and standard deviation (4.06) only, or the test period leaks in
\item The RNN input shape is (batch, steps, features)
\end{itemize}
\begin{lstlisting}
def windows(a, k=7):
    n = len(a) - k
    x = np.stack([a[i:i + k] for i in range(n)])
    return torch.tensor(x).unsqueeze(-1), torch.tensor(a[k:])

>>> tuple(xtr.shape), tuple(xte.shape)
((2913, 7, 1), (730, 7, 1))
\end{lstlisting}
\end{frame}

%%%%%%%%%%%%%%%%%%%%%%%%%%%%%%%%%%%%%%%%%%%%%%%%%%%%%%%%%%
\begin{frame}[fragile]\frametitle{Baselines to Beat}
\begin{itemize}
\item A model is only worth its complexity if it beats the simple answers
\item Error is measured on the 730 test days: mean absolute error (MAE) and root mean squared error (RMSE), in degrees C
\item The linear fit is ordinary least squares on the 7 values and a constant
\end{itemize}
\begin{center}
\small
\begin{tabular}{lcc}
\hline
Method & MAE & RMSE \\
\hline
Same as yesterday & 1.95 & 2.48 \\
Mean of the last 7 days & 2.04 & 2.60 \\
Linear fit on the last 7 days & 1.76 & 2.24 \\
\hline
\end{tabular}
\end{center}
\end{frame}

%%%%%%%%%%%%%%%%%%%%%%%%%%%%%%%%%%%%%%%%%%%%%%%%%%%%%%%%%%
\begin{frame}[fragile]\frametitle{A GRU and an LSTM}
\begin{itemize}
\item A recurrent layer with 16 hidden units reads the 7 days in order
\item Only its output after the last day goes to a single-number linear head
\item GRU: 929 parameters. LSTM: 1,233. Both are small, well under the CNN's 5,258 on digits
\end{itemize}
\begin{lstlisting}
class Net(nn.Module):
    def __init__(self, kind, hidden=16):
        super().__init__()
        cell = nn.GRU if kind == 'GRU' else nn.LSTM
        self.rnn = cell(1, hidden, batch_first=True)
        self.head = nn.Linear(hidden, 1)

    def forward(self, x):          # x: (batch, 7, 1)
        out, _ = self.rnn(x)       # (batch, 7, 16)
        return self.head(out[:, -1]).squeeze(-1)
\end{lstlisting}
\end{frame}

%%%%%%%%%%%%%%%%%%%%%%%%%%%%%%%%%%%%%%%%%%%%%%%%%%%%%%%%%%
\begin{frame}[fragile]\frametitle{Training}
\begin{itemize}
\item Mean squared error loss, Adam with learning rate 0.005, batches of 64, 30 passes over the training windows
\item Most of the learning happens early; after about 5 to 10 epochs the test error stops improving
\end{itemize}
\begin{center}
\small
\begin{tabular}{lccccc}
\hline
Test MAE after epoch & 1 & 5 & 10 & 20 & 30 \\
\hline
GRU & 1.88 & 1.74 & 1.74 & 1.75 & 1.75 \\
LSTM & 1.97 & 1.78 & 1.74 & 1.75 & 1.74 \\
\hline
\end{tabular}
\end{center}
\end{frame}

%%%%%%%%%%%%%%%%%%%%%%%%%%%%%%%%%%%%%%%%%%%%%%%%%%%%%%%%%%
\begin{frame}[fragile]\frametitle{The Result}
\begin{itemize}
\item The recurrent nets are about 11\% better than ``same as yesterday'' (MAE 1.74 against 1.95)
\item They are no better than a straight-line fit on the same 7 days (1.74 and 1.75 against 1.76)
\item Within 1 degree C of the truth: LSTM 36.0\% of test days, same as yesterday 33.8\%
\item The largest single-day miss of the LSTM is 7.1 degrees C
\end{itemize}
\begin{center}
\small
\begin{tabular}{lcc}
\hline
Method & MAE & RMSE \\
\hline
Same as yesterday & 1.95 & 2.48 \\
Mean of the last 7 days & 2.04 & 2.60 \\
Linear fit on the last 7 days & 1.76 & 2.24 \\
GRU & 1.75 & 2.22 \\
LSTM & 1.74 & 2.21 \\
\hline
\end{tabular}
\end{center}
\end{frame}

%%%%%%%%%%%%%%%%%%%%%%%%%%%%%%%%%%%%%%%%%%%%%%%%%%%%%%%%%%
\begin{frame}[fragile]\frametitle{The First 90 Test Days}
\begin{itemize}
\item ``Same as yesterday'' is the actual curve shifted one day to the right: it reaches every peak and dip a day late
\item The LSTM is smoother and leans toward the recent average; it misses the tallest peaks, for example near day 60
\end{itemize}
\begin{center}
\includegraphics[width=0.95\linewidth]{dlcoep_rnn_forecast.png}
\end{center}
\end{frame}

%%%%%%%%%%%%%%%%%%%%%%%%%%%%%%%%%%%%%%%%%%%%%%%%%%%%%%%%%%
\begin{frame}[fragile]\frametitle{What This Tells Us}
\begin{itemize}
\item Here a straight line on 7 days does as well as the RNN, so this small RNN finds little extra structure to use
\item An RNN earns its keep when the input is long, varies in length, or has several signals at once
\item This task has none of those
\item Always run the simple baselines first; a fancy model that only ties them has not earned its complexity
\item For a machine, is the next reading what you need, or a warning that the signal has changed?
\end{itemize}
\end{frame}

%%%%%%%%%%%%%%%%%%%%%%%%%%%%%%%%%%%%%%%%%%%%%%%%%%%%%%%%%%
\begin{frame}[fragile]\frametitle{Quick Check: The Case Study}
\begin{block}{Think About It}
The LSTM's test MAE was 1.74 degrees C. A linear fit on the same 7 days gave 1.76. What is the sensible conclusion?
\begin{enumerate}
  \item[A.] The LSTM is clearly better, so use it everywhere.
  \item[B.] For this task the two are practically tied, so the simpler linear model is the better choice.
  \item[C.] The linear fit must be overfitting.
  \item[D.] The LSTM needs more epochs to win.
\end{enumerate}
\end{block}
\end{frame}

%%%%%%%%%%%%%%%%%%%%%%%%%%%%%%%%%%%%%%%%%%%%%%%%%%%%%%%%%%
\begin{frame}[fragile]\frametitle{Quick Check: The Case Study}
\begin{block}{Answer}
\textbf{Correct: B.} A gap of 0.02 degrees is far below the day-to-day swing of several degrees. The test error stopped improving after about 10 epochs (see the training table), so D does not help, and nothing we measured suggests C.
\end{block}
\end{frame}

%%%%%%%%%%%%%%%%%%%%%%%%%%%%%%%%%%%%%%%%%%%%%%%%%%%%%%%%%%
\begin{frame}[fragile]\frametitle{Try It Yourself}
\begin{itemize}
\item Script: \texttt{Code/\allowbreak dlcoep/\allowbreak session02\_rnn/\allowbreak rnn\_temperature\_case\_study.py}, in the \texttt{dlcoep} conda environment
\item The data file is \texttt{Code/\allowbreak ml/\allowbreak data/\allowbreak daily-minimum-\allowbreak temperatures-in-me.csv}; no download is needed
\item Change one thing and predict the effect first: a 3-day window, a 30-day window, or 4 hidden units
\item Does a longer window help the RNN more than it helps the linear fit?
\end{itemize}
\end{frame}
